% ===========================================================================
%  HOLDING FILE --- NOT COMPILED INTO CHAPTER A
%
%  Excised BDC-facing material from Chapter A (mathematical introduction).
%  Preserved verbatim for possible reuse in later chapters (4a/4b and beyond).
%  Do not \input this file from main.tex or chapter.tex.
%
%  Source audit date: 2026-08-10
%  Editorial horizon: Chapter A orients only for the single-cell BDC chapter
%  (3 BDC_core DRAFT U / main.pdf). Material below anticipates later work:
%  killed-generator viewpoint, mean-field hazard warning, BDC quasi-stationarity,
%  and rupture-into-shared-medium / reseeding.
% ===========================================================================

% ---------------------------------------------------------------------------
% FORMER LOCATION: sections/02_markov_chains.tex
%   Immediately after the birth--death vs BDC rate diagram (m:fig:ratediagram),
%   within subsubsection m:sec:bdc.
% LIKELY DESTINATION: later BDC methods / conditioning chapters; killed-process
%   convention differs from the H-state definition retained in Chapter A.
% ---------------------------------------------------------------------------

%% BEGIN: killed generator, survival-mass balance, and multi-chapter handoff
% (Original notation used catastrophe-to-0 / kappa_i; Chapter A now uses H / delta.)
%
% INTEGRATION NOTE (E2/E3, added 2026-08-10).
%  (a) STOPPING TIME.  T_0 below is the hitting time of 0 under the OLD
%      convention, in which a catastrophe sends the process directly to 0, so
%      T_0 is there already the first of the two mechanisms.  Under Chapter A's
%      convention (m:def:bdc) catastrophe goes to a SEPARATE state H, and every
%      occurrence of T_0 in this block and in the quasi-stationarity block below
%      must become the absorption time  T_0 \wedge T_H  (first of internal
%      extinction and rupture).  Left uncorrected here so the two conventions
%      are not silently mixed; correct it at the point of integration, together
%      with the survival-mass balance, whose right-hand side then splits into a
%      mu_1 term draining into 0 and a catastrophe term draining into H.
%  (b) SYMBOL.  kappa_i must become delta*i on integration.  NOTE that the BDC
%      core chapter (3 BDC_core DRAFT U) already uses kappa for the unrelated
%      reference constant 1 + delta/(2*lambda); kappa is therefore NOT free
%      there and must not survive into the merged text.
Let the positive state $i$ carry
birth rate $\lambda_i$, ordinary death rate $\mu_i$ and catastrophe rate
$\kappa_i$, a catastrophe sending the process directly to $0$. On bounded test
functions the generator is
\begin{equation}
\label{m:eq:bdcgenerator}
(\mathcal Lf)(i)
= \lambda_i\{f(i+1)-f(i)\} + \mu_i\{f(i-1)-f(i)\} + \kappa_i\{f(0)-f(i)\},
\qquad i\ge1,
\end{equation}
with $(\mathcal Lf)(0)=0$. Restricting to the positive states gives the
\emph{killed} subgenerator $K$, with $K_{i,i+1}=\lambda_i$, $K_{i,i-1}=\mu_i$ for
$i\ge2$, and
\begin{equation}
\label{m:eq:killedsubgen}
K_{i,i} = -(\lambda_i+\mu_i+\kappa_i) .
\end{equation}
The row mass missing from $K$ is the instantaneous absorption rate: $\kappa_i$ for
$i\ge2$, and $\mu_1+\kappa_1$ for $i=1$. Writing $T_0=\inf\{t\ge0:X_t=0\}$ and
summing the positive-state forward equations,
\begin{equation}
\label{m:eq:bdcsurvivalloss}
\ddt{}\pr{T_0>t}
= -\mu_1\pr{X_t=1} - \sum_{i\ge1}\kappa_i\pr{X_t=i},
\end{equation}
all probabilities on the right being understood before absorption. Ordinary death
drains survival mass only from state $1$; catastrophe drains it from every positive
state. The
consequences for conditioning are taken up in \cref{m:sec:condmeans}, and the
process itself is developed \fwd{bdc}{in the later chapters on
birth--death--catastrophe processes}, with the two-type extension treated
\fwd{multitype}{in the chapter on multi-type processes} and the release of a
ruptured compartment's contents into a shared medium
\fwd{rupture}{in the chapter on compartment rupture}.
%% END: killed generator block

% ---------------------------------------------------------------------------
% FORMER LOCATION: sections/02_markov_chains.tex, within m:sec:bdc
% LIKELY DESTINATION: discrete-time / generating-function discussion of
%   state-dependent rupture; or a methods warning in a later BDC chapter.
% ---------------------------------------------------------------------------

%% BEGIN: mean-field hazard remark
% Notation aligned with 3 BDC_core DRAFT U / main:
%   X = census of replicative units; \lambda birth, \mu death, \delta catastrophe
%   (per unit); H = absorbing catastrophe/rupture state.
%
% INTEGRATION NOTE (A2, added 2026-08-10).  The version of this remark that is
% compiled in Chapter A no longer reuses \delta for the discrete-time trigger
% probability: it writes \chi \in (0,1) for the per-individual probability per
% generation, reserving \delta for the continuous per-capita rate as in
% m:def:bdc and in the BDC core chapter.  The copy below is the OLD text and
% still overloads \delta; adopt the \chi form when this block is integrated.
\begin{remark}[A state-dependent catastrophe is not a mean-field hazard]
\label{m:rem:notmeanfield}
It is tempting to replace a state-dependent catastrophe rate by a deterministic
hazard evaluated at the mean population, and it does not work. The point is
clearest in discrete time. Write $\delta\in(0,1)$ for the probability that a given
individual initiates catastrophe at a generation --- the discrete counterpart of
the per-capita catastrophe rate $\delta$ of the continuous-time process
\cref{m:def:bdc}, whose total rate out of census $n$ is $\delta n$ into the
absorbing rupture state $H$. Let $\widetilde X_i$ be the branching population with
catastrophe switched off, and let
\begin{equation*}
\mathcal{E}_n = \sum_{i=0}^{n-1}\widetilde X_i
\end{equation*}
be the cumulative exposure through generation $n-1$ (the letter $H$ is reserved
for the catastrophe state itself). Conditioning on the entire catastrophe-free
genealogy, every individual present independently declines to trigger, so
\begin{equation}
\label{m:eq:norupture}
\pr{\text{no catastrophe before }n} = \ex{(1-\delta)^{\mathcal{E}_n}} .
\end{equation}
This identity is exact. Replacing $\mathcal{E}_n$ by $\ex{\mathcal{E}_n}$ is not:
the map $e\mapsto(1-\delta)^e$ is convex, so Jensen's inequality gives
$\ex{(1-\delta)^{\mathcal{E}_n}}\ge(1-\delta)^{\ex{\mathcal{E}_n}}$, and the
substitution systematically understates the chance of getting through: the
expectation in \cref{m:eq:norupture} is a nonlinear functional of the whole path,
and no deterministic hazard closes it.
Two further cautions belong with it. Non-catastrophe is not the same event as
population survival: the probability of being both unruptured and non-zero at
generation $n$ is $\ex{\ind\{\widetilde X_n>0\}(1-\delta)^{\mathcal{E}_n}}$, for
which Jensen supplies no bound at all. And once a trigger occurs every surviving
lineage is stopped together --- the compartment enters $H$ --- so the killed
lineages are coupled and cannot be treated as independent after the fact ---
although their family trees and trigger variables do factorise \emph{before} any
trigger, which is what makes \cref{m:eq:norupture} correct. Closing the model
properly therefore requires a marked or multivariate generating function, or the
killed-generator / rate description of \cref{m:def:bdc,m:eq:bdcrates}.
\end{remark}
%% END: mean-field hazard remark

% ---------------------------------------------------------------------------
% FORMER LOCATION: sections/02_markov_chains.tex, m:sec:coupledexample
%   (end of Coupled ODE--CTMC systems section)
% LIKELY DESTINATION: chapter on compartment rupture / shared medium
%   (beyond the single-cell BDC core chapter).
% FIGURE ASSET retained on disk: figures/rupture_sawtooth.pdf
%
% INTEGRATION NOTE (E1/A2, added 2026-08-10).  SYMBOLS ARE INVERTED in this
% block relative to Chapter A and to the BDC core chapter.  Here \delta is the
% per-unit CLEARANCE rate of the medium (m:eq:mediumdecay) and \rho is the
% CATASTROPHE rate; everywhere else \delta is the catastrophe rate.  The figure
% caption below carries the same inversion (\rho = 0.08 catastrophe,
% \delta = 0.30 clearance), and rupture_sawtooth.pdf was generated under it.
% Before integrating: rename clearance to a free letter (\rho, \chi, \beta and
% \omega are all free in the BDC core; \omega is now used in Chapter A app. B
% and \chi in m:rem:notmeanfield), restore \delta to the catastrophe rate, and
% regenerate the figure so its labels match.
% ---------------------------------------------------------------------------

%% BEGIN: rupture into a shared medium
\subsubsection{A schematic instance: rupture into a shared medium}
\label{m:sec:coupledexample}

It is worth seeing the framework carry one concrete mechanism, and the natural
one here is suggested by \cref{m:def:bdc}. Let $X_t$ be the contents of a
compartment, evolving as the birth--death--catastrophe process of that
definition, and let $y(t)\ge0$ be the concentration of material in a medium the
compartment sits in. Between catastrophes the medium clears at a constant
per-unit rate $\delta$,
\begin{equation}
\label{m:eq:mediumdecay}
\dda{y}{t} = -\delta y ,
\end{equation}
and at a catastrophe time $\tau$ the entire contents of the compartment are
released into the medium at once,
\begin{equation}
\label{m:eq:mediumjump}
y(\tau+) = y(\tau-) + c\,X_{\tau-} ,
\end{equation}
with $c>0$ the material carried per individual. That is the first arrow of the
coupling: the discrete component drives the flow. The second is equally present,
since the catastrophe rate $\rho$, and any rate at which a surviving compartment
takes material back up, may themselves be functions of the current concentration,
so that a fuller medium hastens the next rupture. Both dependences are then
active and the system is of the two-way form \cref{m:eq:coupledthree}.

This is a schematic: it fixes the mechanism and the shape of the state, and the
models actually built on it are developed \fwd{rupture}{in the chapter on
compartment rupture}. Two of its features are worth isolating even so, because
they survive into every instance and neither is visible in
\cref{m:eq:coupledgenerator} as written. The continuous variable is
\emph{impulsive}: \cref{m:eq:mediumjump} is a jump map, not a vector field, so
the natural description is a flow punctuated by resets rather than a single
differential equation, and the generator acquires a corresponding jump term in
$y$. And the quantity released, $X_{\tau-}$, is the state of the chain at a
stopping time rather than a parameter of the model: what the rupture finds is
whatever the conditioned process happens to hold, so its law is the
quasi-stationary problem of \cref{m:sec:condmeans} evaluated at a random time.
That is why systems of this kind are approached with the generating-function and
first-step machinery of \cref{m:sec:methods} rather than closed by simulation
alone. \Cref{m:fig:sawtooth} shows one realisation.

\begin{figure}[ht]
\centering
\includegraphics[width=0.86\textwidth]{rupture_sawtooth.pdf}
\caption{One exact event-driven realisation of the schematic
\cref{m:eq:mediumdecay,m:eq:mediumjump}, with $\lambda=0.55$, $\mu=0.25$,
$\rho=0.08$, $X_0=5$, $\delta=0.30$ and $c=0.4$; so that the picture continues
past the first rupture, the compartment is re-seeded with $X_0$ individuals after
each one. Upper panel: the compartment contents $X_t$, growing between ruptures
and collapsing at them. Lower panel: the medium $y(t)$, decaying exponentially
between ruptures and receiving the impulsive release $c\,X_{\tau-}$ at each.
Dotted verticals mark the rupture times, aligned across the panels. The releases
differ in size because each is the compartment's state at the moment its
catastrophe clock happened to fire, and a rate proportional to the contents
biases that state upwards, a fuller compartment rupturing sooner. That is why
\cref{m:eq:mediumjump} cannot be replaced by a fixed increment.}
\label{m:fig:sawtooth}
\end{figure}
%% END: rupture into a shared medium

% ---------------------------------------------------------------------------
% FORMER LOCATION: sections/03_methods.tex, end of m:sec:condmeans
%   paragraph "Conditioning under two absorbing mechanisms"
% LIKELY DESTINATION: BDC quasi-stationarity discussion after the core
%   single-cell results; depends on killed-subgenerator labels above.
% ---------------------------------------------------------------------------

%% BEGIN: BDC quasi-stationarity paragraph
\paragraph{Conditioning under two absorbing mechanisms.}
One structural point about the birth--death--catastrophe process of
\cref{m:def:bdc} belongs here. If a distribution $\nu$ on the positive states is
quasi-stationary in the sense of \cref{m:def:qsd} and the relevant sums are finite,
then for some $\theta>0$
\begin{equation}
\label{m:eq:bdcqsdeigen}
\nu K = -\theta\nu,
\qquad
\mathbb{P}_\nu\!\left(T_0>t\right) = \ee^{-\theta t},
\end{equation}
with $K$ the killed subgenerator \cref{m:eq:killedsubgen}. Indeed
quasi-stationarity writes the unnormalised positive-state law at time $t$ as
$a(t)\nu$; the semigroup property forces $a(s+t)=a(s)a(t)$, and right-continuity
gives $a(t)=\ee^{-\theta t}$. Differentiating at zero yields the left-eigenvector
relation, and summing its coordinates against \cref{m:eq:bdcsurvivalloss}
identifies the decay rate as
\begin{equation}
\label{m:eq:bdcdecayrate}
\theta = \mu_1\nu_1 + \sum_{i\ge1}\kappa_i\nu_i .
\end{equation}
This is the continuous-time statement of what $A(p)$ and $(2p)^n$ do together in
\cref{m:eq:Apdef}. One limiting case is worth
isolating: if $\kappa_i\equiv\kappa$ is constant, the catastrophe clock is
independent of the birth--death process, the killed semigroup factorises as
$P^K_t = \ee^{-\kappa t}P^{\mathrm{BD},K}_t$, and the extra clock shifts the decay
rate by $\kappa$ while leaving the conditional law unchanged. A state-dependent
sequence such as $\kappa_i=i\kappa$ weights paths by their accumulated exposure and
generally changes the quasi-stationary law itself. Whether a quasi-stationary
distribution survives when the process is conditioned on \emph{neither} extinction
nor rupture is not settled by the argument above, and is best approached by
simulation in the first instance.
%% END: BDC quasi-stationarity paragraph

% ---------------------------------------------------------------------------
% FORMER CROSS-REFERENCE SNIPPET: logistic surrogate paragraph in
%   sections/02_markov_chains.tex (m:sec:logisticspeciation), which cited
%   m:rem:notmeanfield. Rewritten in the compiled chapter without that cite.
% ---------------------------------------------------------------------------
%
% That replacement is a modelling step and not an identity. It is the same
% substitution of a mean for a random census that \cref{m:rem:notmeanfield} shows to
% fail outright for a state-dependent catastrophe, and what separates the two cases is
% only what is being asked of it: here it defines a tractable comparison process, there
% it was asked to return an exact survival probability and did not. Everything that
% follows is accordingly a statement about the surrogate, not about the chain
% \cref{m:eq:speciationrate} defines.
